\chapter{Activity State}

\noindent\textbf{Activity State} estimates \textit{the type and intensity of whole-body
movement} --- sedentary, walking, running, or cycling --- from the strength and frequency
of acceleration. It is theoretically distinct from \textbf{Postural State} (which measures
body orientation, not whether or how the body is moving) and from \textbf{Movement
Regularity} (which measures how rhythmic a motion is, not what kind it is or how energetic).
This is the foundational axis: it labels \emph{what the body is doing}, and the other axes
either build on that label or resolve the questions it deliberately leaves open.

\begin{callout}
\textbf{Scope.} This axis reports the dominant activity \emph{class} over a short window,
not a continuous intensity. Standing folds into \emph{sedentary} on purpose --- posture is a
separate question, delegated to Postural State. Cycling, by contrast, is kept as its own
class because the body is working hard while producing little linear acceleration: without a
dedicated (gyroscope) cue it would read as sedentary, and its raised heart rate would then be
misattributed to stress rather than exercise. The confound most likely to occur is
\textbf{walking versus running}: their energy distributions genuinely overlap, and the split
in that overlap rests on cadence and on placement-specific thresholds. Cycling can only be
detected when a gyroscope is present; without one it reads as sedentary.
\end{callout}

\section{Theoretical Basis}

The rule is an \emph{accelerometer intensity cut-point} classifier, the decades-old method
of the physical-activity-monitoring literature, extended with a cadence tie-breaker and a
rotational cue for cycling. The core feature --- the standard deviation of the acceleration
magnitude --- is the same construct as \textbf{MAD} (Mean Amplitude Deviation,
V\"ah\"a-Yp\"y\"a et al., 2015) and \textbf{ENMO} (Euclidean Norm Minus One, Hildebrand et
al., 2014, the GGIR default), and the cut-point idea traces to activity-count thresholds
(Freedson et al., 1998; Troiano et al., 2008). What is \emph{directly evidenced} by that
work is the use of a magnitude-intensity statistic against fixed thresholds. The
\emph{engineering adaptations}, flagged below, are the specific threshold values (re-derived
on MotionSense), the cadence tie-breaker inside the overlap band, and the gyroscope cycling
branch.

\begin{itemize}[leftmargin=1.4em]
  \item \textbf{Acceleration magnitude intensity} --- $\std(|a_\text{body}|)$ separates
  still from moving and low- from high-energy gait; it is orientation-invariant, so it does
  not depend on how the device is carried (Yurtman \& Barshan, 2017).
  \item \textbf{Cadence (step frequency)} --- where energy alone cannot separate a fast walk
  from a slow run, step rate can; walking and running occupy different cadence ranges in the
  gait literature.
  \item \textbf{Rotational energy for cycling} --- pedalling produces little linear
  acceleration but strong limb/segment rotation, so mean gyroscope magnitude detects cycling
  that the accelerometer misses.
\end{itemize}

\section{Input Signals}

\begin{center}
\begin{tabular}{@{}l l >{\raggedright\arraybackslash}p{7.2cm}@{}}
\toprule
\textbf{Signal} & \textbf{Symbol} & \textbf{Description} \\
\midrule
Body acceleration     & $a_\text{body}$ & Gravity-removed 3-axis acceleration (m/s$^2$) \\
Motion size           & \texttt{move}   & $\std(|a_\text{body}|/g)$ over the window (g) \\
Energy                & \texttt{energy} & $\overline{|a_\text{body}|/g}$, mean magnitude (g) \\
Cadence               & \texttt{cad}    & Step frequency (Hz) from the autocorrelation peak \\
Gyroscope energy      & \texttt{gyro}   & $\overline{|\omega|}$, mean angular-rate magnitude (rad/s); cycling cue \\
\bottomrule
\end{tabular}
\end{center}

\noindent Window $128$ samples ($2.56$\,s at $F_s = 50$\,Hz), 50\% overlap --- the shortest
of the axes, so activity reacts quickly. No personal baseline: the thresholds are
population-derived and frozen.

\section{Rule Formula}

\begin{rulebox}{RULE-AS}{Activity State Classifier}
\textbf{Step 1 --- Features} ($s = |a_\text{body}|/g$; cadence over lag band $[10,64]$):
\[ \texttt{move}=\std(s),\quad \texttt{energy}=\overline{s},\quad
   \texttt{cad}=\frac{F_s}{k^\star},\quad \texttt{gyro}=\overline{|\omega|}. \]

\textbf{Step 2 --- Still gate:}
\[ \texttt{move} < 0.10\text{ g} \ \Rightarrow\ \textbf{sedentary}. \]

\textbf{Step 3 --- Walk / run split} (else, i.e.\ moving):
\[
\begin{aligned}
0.77 \leq \texttt{energy} \leq 1.07:\quad
   &\textbf{running if } \texttt{cad}\geq 1.20\text{ Hz, else } \textbf{walking}\\
\text{otherwise}:\quad
   &\textbf{running if } \texttt{energy}\geq 1.00\text{ g, else } \textbf{walking}
\end{aligned}
\]

\textbf{Step 4 --- Cycling override} (only in the still region, needs gyro):
\[ \big(\texttt{move} < 0.10 \ \wedge\ \texttt{gyro} > 0.18\text{ rad/s}\big)\ \Rightarrow\ \textbf{cycling}. \]

\boxform{ \texttt{label} \in \{\text{sedentary},\ \text{walking},\ \text{running},\ \text{cycling}\} }

\smallskip
\textit{Example} (real window, \S\,1.8): $\texttt{move}=0.54$, $\texttt{energy}=0.91$ (in
the overlap band $\Rightarrow$ energy cannot decide), $\texttt{cad}=1.43\geq1.20\Rightarrow$
\textbf{running}.
\end{rulebox}

\section{Parameter Notes}

\noindent\textbf{Still threshold $T_\text{STILL}=0.10$\,g} {\footnotesize\textsc{[fitted, MotionSense]}}\textbf{.}
Below this, magnitude variation is sensor noise, not body motion. The UCI-HAR value
($0.05$) did \emph{not} transfer because of the unit/placement difference, which is why every
threshold here is stated in g on the same body-acceleration scale.

\noindent\textbf{Run energy line $T_\text{RUN}=1.00$\,g and overlap band $[0.77,\,1.07]$}
{\footnotesize\textsc{[fitted / data, MotionSense]}}\textbf{.}
Outside the band, mean magnitude alone separates walking from running. Inside the band the
two classes' energy distributions touch, so energy is \emph{not} trusted there --- the
cadence tie-breaker takes over.

\noindent\textbf{Cadence split $CAD_\text{RUN}=1.20$\,Hz} {\footnotesize\textsc{[provisional, MotionSense]}}\textbf{.}
Used \emph{only} inside the overlap band. It is placement-specific: on waist data (HHAR) the
autocorrelation fundamental flips between step and stride, so this value does not transfer
unchanged. In practice cadence is a \emph{narrow refinement, not a pillar}: it fires only in
the overlap band, lifted held-out accuracy by about $1.8$ points ($95.3\to97.1\%$), and is
the most fragile part of the rule --- an item to re-derive on own-phone pocket data.

\noindent\textbf{Gyro cycling threshold $T_\text{GYRO}=0.18$\,rad/s} {\footnotesize\textsc{[provisional, HHAR]}}\textbf{.}
Derived on HHAR (waist); pocket cycling is not yet verified. The branch fires only in the
still region so it cannot corrupt walking/running --- but on held-out data this threshold both
\emph{misses} hard pedalling and \emph{false-fires} on phone-handling while standing
(quantified in \S\,1.5), so it is an open validation item.

\noindent\textbf{Fusion method.}
This is a \emph{priority decision cascade}, not a weighted sum. Energy decides first;
cadence is a \emph{tie-breaker} used only where energy is ambiguous; the gyroscope is an
\emph{override gate} confined to the still region. The ordering encodes which cue is trusted
in which regime, keeping each class's decision traceable to one deciding quantity.

\noindent\emph{Tags} (\textsc{fitted}, \textsc{data}, \textsc{prov}, \textsc{conv},
\textsc{design}) classify where each number comes from; their meanings are defined in the Key
Terms front-matter. Window ($128/64$) and cadence lag band ($[10,64]$) are \textsc{conv}
(standard/physiological choices); the confidence formula below is \textsc{design}.

\section{Confidence}

Confidence is a deterministic \emph{margin proxy} in $[0.5,1.0]$: how far the deciding
quantity sits from the threshold that chose the label.

\noindent\textbf{Step C1 --- Deciding margin $m$} (by the branch that fired):
\[
m =
\begin{cases}
(\texttt{gyro}-0.18)/0.18 & \text{cycling}\\[2pt]
(0.10-\texttt{move})/0.10 & \text{sedentary}\\[2pt]
|\texttt{cad}-1.20|/1.20  & \text{walk/run inside the band}\\[2pt]
(\texttt{energy}-1.00)/1.00 & \text{running (outside band)}\\[2pt]
(1.00-\texttt{energy})/1.00 & \text{walking (outside band)}
\end{cases}
\]

\noindent\textbf{Step C2 --- Confidence:}
\boxform{ \text{conf} = 0.5 + 0.5\,\clamp(m,\,0,\,1) }

\noindent\emph{Confidence ceiling / caveat.} This is a \emph{proxy}, not the trained,
calibrated confidence. On the source bench (MotionSense) the calibrated model reaches
$\mathrm{ECE}=0.024$; the proxy reads $0.5$ at a decision boundary (maximum uncertainty) and
$1.0$ far from it. Off-source it is less calibrated, so it should be read as a relative
distance-to-boundary, not a probability.

\medskip
\noindent\emph{Measured per class} (median confidence, held-out data): sedentary $0.97$,
running $0.73$, walking $0.66$, cycling $1.00$ when it fires. Walking is structurally the
least confident --- it is squeezed between the still line below and the run line above, so it
always sits near a boundary. Cycling is confident \emph{when} the gyro cue trips, but that
confidence measures distance-past-threshold, not correctness: on held-out data the branch
\emph{missed} ${\sim}31\%$ of real cycling (hard pedalling crosses the still gate and reads as
walking) and \emph{false-fired} on ${\sim}16\%$ of standing (phone handling pushes gyro past
the provisional $T_\text{GYRO}$). Both are open items to validate on own-phone pocket data.

\section{Interpretation Scale}

Activity State is categorical; the ``scale'' is the class set and each class's signature.

\begin{center}
\begin{tabular}{@{}l >{\raggedright\arraybackslash}p{4.6cm} >{\raggedright\arraybackslash}p{5.2cm}@{}}
\toprule
\textbf{Class} & \textbf{Signature} & \textbf{Meaning} \\
\midrule
Sedentary & $\texttt{move}<0.10$\,g & Still or near-still: sitting, standing, resting (posture delegated) \\
Walking   & Moving, below the run line / in-band cadence $<1.20$\,Hz & Steady low-to-moderate gait \\
Running   & $\texttt{energy}\geq 1.00$\,g, or in-band cadence $\geq 1.20$\,Hz & High-energy, fast-cadence gait \\
Cycling   & Still-region body motion with $\texttt{gyro}>0.18$\,rad/s & Low linear motion but rotational (pedalling); needs a gyroscope \\
\bottomrule
\end{tabular}
\end{center}

\section{Context Applicability}

\begin{callout}
\textbf{Target context:} a phone or wearable carried on the body (pocket-scale placement),
sampling accelerometer and gyroscope at $50$\,Hz.
\end{callout}

\begin{center}
\begin{tabular}{@{}>{\raggedright\arraybackslash}p{5.4cm} c c >{\raggedright\arraybackslash}p{3.8cm}@{}}
\toprule
\textbf{Context} & \textbf{accel} & \textbf{gyro} & \textbf{Reliable?} \\
\midrule
Phone in pocket (target)      & Yes & Yes     & Yes --- all four classes \\
Waist / no gyroscope          & Yes & No      & Partial --- cycling reads as sedentary \\
Wrist / other placement       & Yes & Partial & Partial --- cadence/gyro thresholds shift \\
\bottomrule
\end{tabular}
\end{center}

\noindent\textbf{Scientific boundary.} Even with ideal data the axis cannot separate
\emph{standing} from \emph{sitting} or from sedentary --- that needs orientation, delegated
to Postural State. The walking/running boundary is a real energy overlap resolved
heuristically by cadence, and both the cadence and gyro thresholds are placement-specific
rather than universal. Cycling additionally requires a placement that is rotationally active
but linearly quiet during pedalling: the torso qualifies (chest), the limbs do not
(wrist/ankle keep moving and read as walking), so cycling detection is torso-placement-specific.

\section{Worked Example: Resolving a Jog in the Overlap Band}

\noindent\textit{Scenario (real captured window).} A user jogs at a steady pace, phone in the
front pocket. The window below is a genuine capture, not an invented one --- MotionSense
subject~19 (held out from all threshold derivation), jog trial, $58.9$\,s in
(\texttt{jog\_9/sub\_19.csv}, samples $2944$--$3071$) --- run through the live axis. Its
energy lands inside the walk/run overlap band, so energy alone cannot decide the class.

\begin{center}
\begin{tabular}{@{}ll@{}}
\toprule
\textbf{Signal} & \textbf{Value} \\
\midrule
\texttt{move} (std of $s$)   & $0.54$\,g \\
\texttt{energy} (mean of $s$)& $0.91$\,g \;(in band $[0.77,1.07]$) \\
\texttt{cad}                 & $1.43$\,Hz \\
\texttt{gyro}                & $1.87$\,rad/s \\
\bottomrule
\end{tabular}
\end{center}

\noindent\textbf{Step-by-step:}
\begin{align*}
\texttt{move}=0.54 &\geq 0.10 \Rightarrow \text{not sedentary}\\
\texttt{energy}=0.91 &\in[0.77,1.07] \Rightarrow \text{use cadence, not energy}\\
\texttt{cad}=1.43 &\geq 1.20 \Rightarrow \textbf{running}\\
\texttt{move} &< 0.10 \text{ false} \Rightarrow \text{cycling override does not fire}\\
m = |1.43-1.20|/1.20 = 0.19,\quad
\text{conf} &= 0.5 + 0.5\cdot 0.19 = \mathbf{0.60}
\end{align*}

\begin{callout}
\textbf{Interpretation:} the window is classified \emph{running}. Because the energy fell in
the overlap band, the decisive driver was the \emph{cadence} ($1.43$\,Hz), not the energy.
Confidence is moderate ($0.60$): the step rate is above but not far above the $1.20$\,Hz
walk/run line, which is exactly the region where this axis is least certain. The gyroscope
reads high ($1.87$\,rad/s) but is ignored --- the cycling override only fires in the still
region, and here the body is clearly moving.
\end{callout}

\section{Implementation Notes}

\begin{itemize}[leftmargin=1.4em]
  \item \textbf{Window.} $128$ samples ($2.56$\,s) at $50$\,Hz, 50\% overlap --- the shortest
  window in the pack, so the activity label updates fastest. There is an inherent latency: a
  $2.56$\,s window means an activity change is only caught within about $1.3$--$2.6$\,s (a new
  reading every ${\sim}1.3$\,s). This matches the physiology it contextualises (heart rate and
  affect shift over tens of seconds) but is \emph{not} intended for instantaneous gesture or
  fall detection.
  \item \textbf{Fallback behavior.} No accel/linear stream $\rightarrow$ null. No gyroscope
  $\rightarrow$ the cycling branch is disabled and the reading note says so; the accel classes
  are unaffected. Cadence still computes from the accelerometer.
  \item \textbf{Relationship to other chapters.} The standing/sitting ambiguity it folds into
  sedentary is resolved by \textbf{Postural State}; its walking/running output gates
  \textbf{Movement Regularity} (only meaningful while moving) and feeds \textbf{Locomotion}
  and the \textbf{Overall State} combiner.
  \item \textbf{Validation status.} Headline \textbf{97.1\%} on MotionSense (phone in pocket,
  held-out subjects; macro-F1 $0.948$, well-calibrated $\mathrm{ECE}=0.024$). The rule then
  holds \emph{frozen} (nothing re-derived) across two further datasets and sensor types
  (table below) --- notably the walk/run boundary, which HHAR cannot test (it has no running),
  is cross-validated on PAMAP2. Cycling is the one placement-dependent part.
  $CAD_\text{RUN}$ and $T_\text{GYRO}$ stay provisional; pocket cycling and a phone-in-pocket
  walk/run set (e.g.\ WISDM) are the remaining checks.
\end{itemize}

\begin{center}
\begin{tabular}{@{}>{\raggedright\arraybackslash}p{2.3cm} >{\raggedright\arraybackslash}p{3.5cm} >{\raggedright\arraybackslash}p{2.5cm} >{\raggedright\arraybackslash}p{4.1cm}@{}}
\toprule
\textbf{Dataset} & \textbf{Sensor / placement} & \textbf{Classes} & \textbf{Result (frozen)} \\
\midrule
MotionSense & phone, pocket & sed / walk / run & $97.1\%$ held-out (macro-F1 $0.948$) \\
HHAR & phones + watches, waist/hand & sed / walk (no running) & macro-F1 $0.996$ shared \\
PAMAP2 & research IMUs, wrist/chest/ankle & walk / run & $0.89$--$0.96$ acc., all 3 placements \\
PAMAP2 & \quad'' \;(cycling) & cycling & chest recall $0.67$; wrist/ankle ${\sim}0.03$ \\
\bottomrule
\end{tabular}
\end{center}

\noindent{\footnotesize PAMAP2 figures use offline per-segment gravity removal; the live
on-device per-window approximation lowers only ankle-walking ($0.72$ vs $0.93$) --- running
and the torso/wrist placements are unaffected.}

% ---- Formula flow diagram ----
\begin{figure}[h]
\centering
\begin{tikzpicture}[
  node distance=10mm and 26mm,
  box/.style={draw, dashed, rounded corners=1pt, align=center, inner sep=5pt, font=\small},
  op/.style={draw, rounded corners=1pt, align=center, inner sep=5pt, font=\small},
  resultbox/.style={draw, double, rounded corners=3pt, align=center, inner sep=6pt, font=\small\bfseries},
  ar/.style={-{Latex[length=2mm]}}]

  \node[box] (a) {Accel magnitude\\ $s=|a_\text{body}|/g$};
  \node[op, below=of a] (feat) {Features:\\ \texttt{move}, \texttt{energy}, \texttt{cad}};
  \node[op, below=of feat] (casc) {Decision cascade\\ (Steps 2--3)};
  \node[resultbox, below=of casc] (out) {\texttt{label} $\in$ \{sedentary, walking,\\ running, cycling\}};

  \node[box, right=of a] (g) {Gyroscope\\ $|\omega|$};
  \node[op, below=of g] (cyc) {Cycling override\\ still $\wedge\ \texttt{gyro}>0.18$\\ (Step 4)};

  \draw[ar] (a) -- (feat);
  \draw[ar] (feat) -- (casc);
  \draw[ar] (casc) -- (out);
  \draw[ar] (g) -- (cyc);
  \draw[ar] (cyc) |- (casc) node[pos=0.28,above,font=\footnotesize] {override};
\end{tikzpicture}
\caption*{\small Formula flow --- Activity State: a priority decision cascade over
orientation-invariant accelerometer features, with a gyroscope override for cycling.}
\end{figure}
